\section{Discusión}
\label{sec:discusion}
% Responde: ítem 8

\paragraph{Frecuencia cardíaca.}
La frecuencia cardíaca media pasó de \SI{84}{lpm} en reposo a \SI{131}{lpm} después del ejercicio, y durante los tres minutos de recuperación bajó de 146 a \SI{120}{lpm} sin volver al valor de reposo. Esto es lo que predice la teoría (sección~\ref{sec:hrv}). Durante el esfuerzo la frecuencia sube por la retirada vagal y la activación simpática; al terminar, la reactivación vagal la hace bajar rápido al principio, y la actividad simpática y las catecolaminas circulantes, que se eliminan más despacio, la mantienen por encima del reposo durante varios minutos \parencite{imai1994,michael2017}. La caída más rápida se dio en el primer minuto (\SI{-14.5}{lpm/min}), como corresponde a la fase vagal de la recuperación. Dos detalles no se ajustan a un modelo simple de caída exponencial: la meseta de unos \SI{20}{\second} al comienzo del registro y la pendiente casi nula del segundo minuto. Una explicación posible para la meseta es que el registro empezó apenas terminado el esfuerzo, cuando la frecuencia todavía no había empezado a bajar, o que el voluntario se estaba moviendo y acomodando en ese momento. No registramos cuánto tiempo pasó entre el fin del ejercicio y el comienzo de la grabación, ni qué hizo el voluntario durante la recuperación, así que no podemos confirmarlo. Con un solo registro tampoco podemos distinguir si el estancamiento del segundo minuto es una característica de la recuperación o una fluctuación.

\paragraph{Variabilidad en el dominio del tiempo.}
La variabilidad latido a latido cayó fuertemente después del ejercicio: el RMSSD bajó de 48 a \SI{8.6}{\milli\second} y el pNN50, de \SI{23}{\percent} a cero. Como ambos índices reflejan sobre todo la modulación vagal \parencite{shaffer2017}, esta caída es consistente con un tono vagal muy reducido en la recuperación temprana. Parte de la reducción se debe a la interacción entre la frecuencia cardíaca y su variabilidad: con intervalos más cortos, una misma modulación produce variaciones absolutas menores \parencite{sacha2014}. Sin embargo, ese efecto no alcanza para explicarla: aun relativa al RR medio, la variabilidad sin tendencia cayó del \SI{11.5}{\percent} al \SI{2.6}{\percent}. El SDNN del registro post ejercicio es engañoso si se lo toma solo, porque está dominado por la tendencia de la recuperación; quitada la tendencia, la reducción respecto del reposo es de siete veces.

\paragraph{Variabilidad en el dominio de la frecuencia (ítem 7a).}
Las potencias LF y HF se redujeron entre 50 y 100 veces después del ejercicio. El espectro de reposo tiene un pico definido en la banda HF, en \SI{0.18}{\hertz}, que corresponde a unas 11 respiraciones por minuto, un valor normal en reposo. Su presencia es la firma de la arritmia sinusal respiratoria y, por lo tanto, de un control vagal activo. En el primer minuto post ejercicio la banda HF prácticamente desapareció (\SI{7}{\percent} de LF + HF) y su pico se corrió a \SI{0.36}{\hertz}, unas 22 respiraciones por minuto. Ambas cosas coinciden con lo esperado: el vago está casi retirado, y la poca modulación respiratoria que queda sigue a una respiración agitada. En los minutos 2 y 3 la potencia HF reapareció y su pico volvió a \SIrange{0.19}{0.21}{\hertz}, lo que indica la reactivación vagal gradual mientras la frecuencia cardíaca baja \parencite{michael2017}. El aumento del cociente LF/HF en el primer minuto (13) es coherente con este cuadro, pero no debe leerse como una medida directa del predominio simpático: refleja sobre todo la caída de HF, y el propio cociente es un indicador discutido del balance autonómico \parencite{billman2013}.

El resultado que menos se ajusta a lo esperado es el cociente LF/HF del reposo, de 4,7. En reposo suele encontrarse un cociente cercano a 1 o 2, con una banda HF importante. En nuestro registro, la banda LF está dominada por tres oscilaciones lentas grandes, de unos \SI{16}{\second} de período, que se ven directamente en el tacograma. Pueden deberse a respiraciones profundas o irregulares, a que el voluntario habló o se movió durante el registro, o a que estaba recuperándose de alguna actividad previa. Además, el registro dura \SI{55}{\second}, menos que los \SI{2}{\minute} que la \textcite{taskforce1996} pide para estimar LF, de modo que la estimación de esa banda descansa en apenas tres ciclos y es poco confiable. Por eso, en la comparación entre condiciones damos más peso a las potencias absolutas de HF, al RMSSD y a la evolución minuto a minuto que al cociente LF/HF del reposo.

\paragraph{Comparación entre voluntarios (ítems 6b y 7b).}
La consigna pide comparar dos voluntarios, pero se registró solo uno. No podemos decir, entonces, si el comportamiento observado es general o propio de V1. Sí podemos señalar qué se esperaría: el patrón cualitativo (aumento de la frecuencia, caída de la VFC, desaparición y recuperación de HF) debería repetirse, mientras que los valores absolutos dependen mucho de cada sujeto. Por ejemplo, en personas entrenadas la reactivación vagal después del ejercicio es más rápida \parencite{imai1994}.

\paragraph{Limitaciones.}
Además de contar con un solo voluntario, el análisis tiene otras limitaciones. Los registros son cortos: \SI{1}{\minute} en reposo y \SI{3}{\minute} post ejercicio, por debajo de los \SI{5}{\minute} que la \textcite{taskforce1996} recomienda como estándar para registros cortos. La frecuencia de muestreo efectiva no es exactamente \SI{500}{\hertz}, porque entre paquetes se pierden muestras: esto comprime el eje de tiempo un \SI{4}{\percent} y agrega una incertidumbre de hasta \SI{8}{\milli\second} en cada intervalo (sección~\ref{sec:fs}). El error es sistemático y no cambia la comparación entre condiciones, pero sí los valores absolutos de frecuencia cardíaca y de las frecuencias del espectro, y la incertidumbre por intervalo es del orden del propio RMSSD post ejercicio, lo que vuelve menos preciso ese valor. El registro post ejercicio no es estacionario, lo que viola uno de los supuestos del análisis espectral; lo mitigamos quitando la tendencia y analizando ventanas de un minuto, que a su vez tienen poca resolución en frecuencia ($1/\SI{60}{\second} \approx \SI{0.017}{\hertz}$). Por último, no se registró la respiración, de modo que la asociación entre el pico HF y la frecuencia respiratoria es una inferencia.

\paragraph{¿Se corresponden los resultados con la teoría?}
En lo esencial, sí. La frecuencia cardíaca aumentó con el ejercicio y bajó durante la recuperación, más rápido al principio; la variabilidad, en especial la ligada al vago (RMSSD, pNN50 y potencia HF), cayó fuertemente después del esfuerzo y empezó a recuperarse en el transcurso de los tres minutos, y el pico respiratorio se desplazó hacia frecuencias más altas cuando la respiración estaba agitada. La principal discrepancia, el cociente LF/HF elevado en reposo, se explica razonablemente por la corta duración del registro y por las oscilaciones lentas que contiene, más que por un fallo del procesamiento.
